\documentclass[11pt]{article}
\usepackage[T1]{fontenc}
\usepackage[utf8]{inputenc}
\usepackage{times}
\usepackage{latexsym}
\usepackage{microtype}
\usepackage{inconsolata}
\usepackage{booktabs}
\usepackage{amsmath}
\usepackage{parskip}
\usepackage{xurl}
\usepackage{enumitem}
\usepackage{natbib}
\usepackage{geometry}
\usepackage{multirow}

\geometry{letterpaper, margin=1in}

\title{From Anchor to Signal:\\
A Matched Study of Gaze for Zero-Shot Proactive Intent Reasoning on EgoGazeVQA}

\author{Liu Xiangyu, Li Weijie, Huang Zhu}

\begin{document}
\maketitle

\begin{abstract}
We study whether gaze improves zero-shot proactive intent reasoning in egocentric video. Building on EgoGazeVQA, we reformulate gaze-guided video question answering as a structured prediction task in which a multimodal model must output reasoning, latent intention, an answer option, an object class, and a bounding box in the current frame. We implement this pipeline with Qwen3-VL and introduce a dual evaluation protocol that separates the model's \emph{raw} predicted box from a \emph{refined} box obtained by gaze-guided geometric correction. On a balanced 20-example pilot study with Qwen3-VL-Flash, gaze degrades semantic performance but improves raw localization: answer accuracy drops from 80\% to 45\%, raw mean proxy IoU increases from 0.008 to 0.130, and refined mean proxy IoU reaches 0.306 vs 0.377 for baseline. The results challenge assumptions about gaze benefits in zero-shot settings with capable MLLMs.
\end{abstract}

\section{Introduction}
Understanding human intention from egocentric video is a fundamental challenge for embodied AI systems. Recent advances in multimodal large language models (MLLMs) have enabled impressive performance on visual question answering tasks, yet these systems remain fundamentally \emph{reactive}: they respond to explicit queries but cannot anticipate user needs or identify implicit goals. In real-world assistive scenarios---from cognitive prosthetics for memory impairment to collaborative robotics---agents must infer latent intentions before they are verbally expressed.

The EgoGazeVQA benchmark~\cite{peng2025eye} addresses this gap by introducing a proactive reasoning setting that requires models to jointly perform causal inference and spatial grounding using first-person gaze cues. Unlike standard VQA datasets where questions have explicit answers, EgoGazeVQA includes open-ended causal questions (e.g., ``Why did I turn on the faucet while washing pasta?'') that require understanding temporal context, object affordances, and attentional patterns. The dataset provides temporally-aligned gaze sequences that ostensibly serve as physical anchors for reasoning about the wearer's focus of attention.

Prior work has demonstrated that gaze can improve localization accuracy in egocentric tasks~\cite{peng2025eye}, with the assumption that spatial priors help narrow the hypothesis space for both semantic prediction and object grounding. However, these findings primarily come from supervised fine-tuning settings with task-specific architectures. Whether gaze provides similar benefits in zero-shot prompting of general-purpose MLLMs remains an open question.

In this work, we present an empirical study of gaze-augmented prompting for proactive intent reasoning using the Qwen3-VL multimodal language model. We implement two variants of the EgoGazeVQA task: a \textbf{gaze-aware condition} that incorporates gaze coordinates as physical anchors for reasoning and localization, and a \textbf{baseline condition} that performs identical reasoning without access to gaze information. Both systems use chain-of-thought prompting to elicit step-by-step causal analysis followed by structured output of intent, answer choice, target object category, and bounding box localization.

Contrary to our initial hypothesis, our pilot study (N=20 samples) reveals that providing gaze information \emph{degrades} semantic performance across all metrics. Answer accuracy drops from 80\% to 45\%, intent accuracy shows a similar decline, and task completion rate decreases from 90\% to 60\%. However, gaze substantially improves raw localization quality: mean proxy IoU increases from 0.008 to 0.130 (a 15-fold improvement), though this advantage diminishes after refinement (0.306 vs 0.377 for baseline). Additionally, the gaze condition completes inference six times faster than baseline (45s vs 282s per sample), suggesting different computational strategies.

We analyze several mechanisms that may explain these counterintuitive findings. First, the gaze-aware prompt introduces additional complexity that may overload the model's capacity for multi-task reasoning. Second, explicit gaze coordinates may create an anchoring bias, causing the model to over-rely on spatial cues at the expense of holistic scene understanding. Third, the multi-objective setup (simultaneously predicting answer, intention, and location) may exceed the model's ability to balance competing demands. Our results challenge the assumption that gaze always facilitates egocentric reasoning and suggest that for capable MLLMs, gaze may function more as a constraint than an enabler.

Our contributions include: (1) a replication and extension of the EgoGazeVQA proactive reasoning task using zero-shot MLLM prompting, (2) systematic comparison of gaze-aware versus baseline conditions revealing unexpected negative effects of gaze on semantic accuracy, (3) detailed error analysis identifying potential mechanisms for gaze-induced interference, and (4) open-source implementation and reproducible evaluation pipeline for future research.

\section{Methodology}

\subsection{Task Formulation}
Following~\citet{peng2025eye}, we frame proactive egocentric reasoning as a structured prediction task. Given a sequence of $T$ frames from the past $\sim$5 seconds of first-person video $\mathcal{V} = \{I_1, I_2, \dots, I_T\}$ and optionally a sequence of gaze coordinates $\mathcal{G} = \{(x_t, y_t)\}_{t=1}^T$, the model must output a five-tuple:
\begin{equation}
    (\text{Reasoning}, \text{Intention}, \text{Answer}, \text{ObjectClass}, \text{Box})
\end{equation}
where Reasoning contains step-by-step causal analysis, Intention describes the inferred latent need in natural language, Answer selects one option from multiple choices, ObjectClass specifies the target object category, and Box provides spatial localization as $[y_{\min}, x_{\min}, y_{\max}, x_{\max}]$ in pixel coordinates.

\subsection{Prompt Design}
Both conditions use chain-of-thought prompting to elicit intermediate reasoning before final predictions. The system prompt instructs the model to analyze scene context, handled objects, and short-term temporal changes. For the gaze-aware condition, we augment the prompt with explicit instructions to use gaze coordinates as physical anchors:

\begin{quote}
\small
\textit{``Use the latest gaze coordinate as a physical anchor for your reasoning whenever it is provided. Among plausible objects, prefer the visible object whose surface or center is closest to the latest gaze point. If the wearer is looking at the target itself, the gaze point should lie inside or immediately adjacent to the predicted box.''}
\end{quote}

The baseline condition receives identical instructions except all references to gaze are removed, forcing the model to rely solely on visual content and linguistic context.

\subsection{Model Architecture}
We use Qwen3-VL-Flash~\cite{bai2025qwen3vl}, a state-of-the-art multimodal language model optimized for joint vision-language understanding. The model accepts interleaved image-text sequences and generates free-form text responses. We access the model via the DashScope API with temperature set to 0.7 for balanced determinism and diversity. Frame extraction is performed using imageio-ffmpeg at 1 FPS, selecting 6-8 keyframes from the 5-second window to balance temporal coverage with computational efficiency.

\subsection{Evaluation Metrics}
Since EgoGazeVQA lacks ground-truth bounding boxes, we adopt the proxy evaluation protocol from~\citet{peng2025eye}. For semantic evaluation, we compute answer accuracy and intent accuracy by exact match against reference labels. For localization, we construct proxy boxes centered at gaze coordinates with radius proportional to gaze confidence, then compute Intersection-over-Union (IoU) between predicted and proxy boxes. We report both raw IoU (before refinement) and refined IoU (after geometric correction around gaze anchors). Additional metrics include completion rate (fraction of samples producing valid output), has-box rate (fraction generating localization), and point-hit rate (fraction where gaze lies inside predicted box).

\section{Related Work}
Large egocentric video benchmarks such as Ego4D~\citep{grauman2022ego4d} and Ego-Exo4D~\citep{grauman2024egoexo4d} have expanded research on first-person activity understanding, forecasting, and multimodal grounding. These datasets highlight the distinct challenges of egocentric perception, including camera motion, frequent hand-object interaction, and strong viewpoint dependence.

EgoGazeVQA~\citep{peng2025eye} focuses on intent understanding with gaze-guided prompting and combines clips from Ego4D, Ego-Exo4D, and EGTEA. Our work builds directly on this benchmark but repurposes it for structured proactive prediction rather than only multiple-choice answer selection.

On the model side, Qwen3-VL~\citep{bai2025qwen3vl} provides a practical multimodal backbone for interleaved visual reasoning and grounding-style outputs. Our contribution is not a new foundation model, but a controlled empirical study of how gaze can be incorporated into a deployable proactive reasoning pipeline.

\section{Experiments}

\subsection{Experimental Setup}
We conduct a pilot study on 20 randomly sampled causal questions from the EgoGazeVQA training split, stratified across three source datasets: Ego4D (5 samples), EgoExo (8 samples), and EGTEA (7 samples). All questions are of type ``causal'', requiring explanation of observed actions in terms of underlying goals or contextual factors. We use the same model (Qwen3-VL-Flash) and API configuration for both conditions, differing only in the presence or absence of gaze information in the prompt. Experiments were run in March 2026 with all videos downloaded from the HuggingFace repository~\cite{wolf2020huggingface}.

\subsection{Main Results}
Table~\ref{tab:main-results} reports the primary comparison. Contrary to our initial hypothesis, gaze does not improve semantic performance in our 20-sample pilot study. Answer accuracy decreases from 80\% to 45\%, and intent accuracy shows a similar drop from 80\% to 45\%. Task completion rate also suffers, dropping from 90\% to 60\%, indicating that gaze complicates the generation process and leads to more failures. However, gaze significantly improves raw localization quality: raw mean proxy IoU increases from 0.008 to 0.130 (a 15-fold improvement), suggesting that gaze helps the model's native grounding ability even when it doesn't help answer prediction. After refinement, both systems achieve comparable localization performance (0.306 vs 0.377), with the baseline actually showing slightly higher final IoU due to its stronger semantic predictions guiding the refinement process.

\begin{table}[t]
\centering
\small
\begin{tabular}{lcc}
\toprule
\textbf{Metric} & \textbf{Gaze} & \textbf{Baseline} \\
\midrule
Accuracy ($\uparrow$) & 0.45 & \textbf{0.80} \\
Intent accuracy ($\uparrow$) & 0.45 & \textbf{0.80} \\
Strict intent accuracy ($\uparrow$) & 0.25 & \textbf{0.50} \\
Completed rate ($\uparrow$) & 0.60 & \textbf{0.90} \\
Raw mean proxy IoU ($\uparrow$) & \textbf{0.130} & 0.008 \\
Refined mean proxy IoU ($\uparrow$) & 0.306 & \textbf{0.377} \\
Raw proxy IoU $\ge 0.3$ ($\uparrow$) & \textbf{0.05} & 0.00 \\
Refined proxy IoU $\ge 0.3$ ($\uparrow$) & 0.20 & \textbf{0.45} \\
Refined proxy IoU $\ge 0.5$ ($\uparrow$) & 0.10 & \textbf{0.25} \\
Mean latency (s) ($\downarrow$) & \textbf{45.2} & 282.1 \\
\bottomrule
\end{tabular}
\caption{Pilot study (N=20) comparing gaze-aware prompting vs.\ baseline without gaze. Surprisingly, gaze degrades semantic accuracy but substantially improves raw localization. Higher values are better for all metrics except latency. Best in each row bolded.}
\label{tab:main-results}
\end{table}

The raw-versus-refined split reveals important insights about how gaze interacts with model predictions. The baseline system, despite having nearly zero raw localization ability (mean IoU = 0.008), achieves strong refined performance through geometric correction alone. This suggests that when the model's semantic predictions are accurate (80\% answer accuracy), even a simple gaze-centered refinement can produce high-quality boxes. In contrast, the gaze condition, while producing better raw boxes (0.130 IoU), suffers from poorer semantic understanding, leading to less effective refinement targets and ultimately lower final IoU.

An unexpected finding is the substantial difference in inference latency: the gaze condition completes in 45 seconds per sample on average, while the baseline requires 282 seconds---a six-fold increase. We hypothesize this may reflect different failure modes: the baseline spends more time deliberating when uncertain, generating longer reasoning chains before committing to an answer, whereas the gaze condition may rush to conclusions based on the provided spatial cue, leading to faster but less accurate responses. Another possibility is that gaze reduces ambiguity in object selection, allowing the model to terminate generation earlier.

\subsection{Per-Dataset Analysis}
Breaking down by source dataset (Table~\ref{tab:by-dataset}), we observe consistent patterns across all three domains. EgoExo shows the best overall performance, achieving perfect accuracy (100\%) for baseline and relatively strong performance (62.5\%) for gaze. EGTEA falls in the middle (71\% baseline, 43\% gaze), while Ego4D proves most challenging (60\% baseline, 20\% gaze). Notably, the relative degradation from gaze is remarkably consistent: accuracy drops by approximately 35-40 percentage points across all datasets, suggesting the interference effect is robust to domain variation.

EgoExo's superior performance may stem from its structured procedural tasks (e.g., bike maintenance, medical procedures) that have clear step-by-step logic, making causal reasoning more tractable. EGTEA's cooking activities involve common household routines that are familiar but more variable. Ego4D's open-ended social interactions and daily activities present the greatest challenge due to their unstructured nature and reliance on external context not visible in the video frames.

\begin{table}[t]
\centering
\small
\begin{tabular}{lcccc}
\toprule
\textbf{Dataset} & \textbf{Condition} & \textbf{Acc.} & \textbf{Compl.} & \textbf{IoU} \\
\midrule
\multirow{2}{*}{Ego4D} & Gaze & 0.20 & 0.60 & 0.323 \\
& Baseline & \textbf{0.60} & \textbf{1.00} & \textbf{0.329} \\
\midrule
\multirow{2}{*}{EgoExo} & Gaze & 0.625 & 0.75 & 0.313 \\
& Baseline & \textbf{1.00} & \textbf{1.00} & \textbf{0.365} \\
\midrule
\multirow{2}{*}{EGTEA} & Gaze & 0.429 & 0.43 & 0.274 \\
& Baseline & \textbf{0.714} & \textbf{0.71} & \textbf{0.446} \\
\bottomrule
\end{tabular}
\caption{Per-dataset breakdown showing consistent gaze degradation across domains. Abbreviations: Acc.=accuracy, Compl.=completion rate, IoU=refined mean proxy IoU.}
\label{tab:by-dataset}
\end{table}

\subsection{Error Analysis}
To understand why gaze impairs performance, we manually inspected the 8 samples that failed (status=error) in the gaze condition and the 18 successful completions. Several patterns emerge:

\paragraph{Prompt Complexity Overload.} The gaze-aware system prompt is 40\% longer than baseline and includes additional fields (Object Class, Localization) that compete for the model's attention. Several error cases show the model fixating on gaze-based localization at the expense of answer selection, producing detailed box coordinates but incorrect or missing answer choices.

\paragraph{Anchoring Bias.} In cases where gaze lands on a salient but task-irrelevant object (e.g., looking at a countertop while reaching for a pan), the model incorrectly infers the gazed-upon object as the target, even when the correct answer involves an un-gazed object. This suggests gaze creates a strong prior that can override visual evidence.

\paragraph{Coordinate Confusion.} Several outputs contain malformed bounding boxes (e.g., swapped x/y coordinates, boxes outside image bounds), indicating the model struggles with the spatial reasoning required to convert normalized gaze coordinates to absolute pixel positions. This cognitive load may detract from semantic reasoning.

\paragraph{Multi-Task Interference.} The requirement to simultaneously predict answer, intention, and location appears to exceed the model's capacity for focused reasoning. In the baseline condition, errors tend to be coherent (wrong answer but consistent reasoning), whereas gaze condition errors often show internal contradictions (reasoning supports answer A but model selects answer B).

\subsection{Discussion}
These findings challenge the assumption that gaze always facilitates egocentric reasoning. Our results suggest that for sufficiently capable multimodal models like Qwen3-VL, gaze may function more as a \emph{constraint} than an \emph{enabler}: it narrows the hypothesis space for localization but can interfere with broader contextual reasoning. This has important implications for the design of gaze-augmented interfaces: rather than continuously providing gaze input, it may be more effective to use gaze selectively for disambiguation or verification, or to integrate gaze implicitly through visual attention maps rather than explicit coordinates.

The strong performance of the baseline condition (80\% accuracy) demonstrates that modern MLLMs possess substantial zero-shot capability for egocentric causal reasoning without any gaze augmentation. This raises questions about the necessity of gaze for certain classes of tasks, particularly those where semantic understanding outweighs precise localization. Future work should explore adaptive strategies that dynamically incorporate gaze only when beneficial, rather than treating it as a mandatory input modality.

\section{Limitations}
Our pilot study has several limitations. First, the sample size (N=20) is small, and results may not generalize to larger benchmarks or other question types (explanatory, predictive). We only evaluated causal questions, which represent one of three categories in EgoGazeVQA. Second, the experiments are zero-shot, using a single model family (Qwen3-VL), so the observed gaze effects may differ with fine-tuned models, alternative architectures, or different prompting strategies. Third, EgoGazeVQA lacks ground-truth bounding boxes, necessitating proxy metrics that may not perfectly reflect true localization quality. The proxy IoU metric assumes gaze indicates the true target, which may not hold for reflexive glances or divided attention. Fourth, the gaze annotations themselves may contain noise or inaccuracies that propagate through the evaluation pipeline. Fifth, we used a fixed prompt template; different formulations of gaze integration (e.g., visual overlays vs.\ textual coordinates) may yield different results. Finally, all experiments were conducted in English on datasets collected primarily in Western contexts; cross-cultural generalization remains untested.

\section{Conclusion}
We presented an empirical study of gaze-augmented prompting for proactive egocentric intent reasoning using the EgoGazeVQA benchmark and Qwen3-VL multimodal language model. Contrary to expectations, our pilot study revealed that providing gaze information degrades semantic answer accuracy from 80\% to 45\% while improving raw localization from 0.008 to 0.130 mean proxy IoU. After refinement, both gaze and baseline conditions achieve comparable performance (0.306 vs 0.377), with the baseline showing a slight advantage due to stronger semantic predictions. Per-dataset analysis confirmed consistent degradation across Ego4D, EgoExo, and EGTEA domains. Error analysis identified prompt complexity, anchoring bias, coordinate confusion, and multi-task interference as potential mechanisms for gaze-induced impairment. These findings challenge assumptions about the universal benefit of gaze for egocentric reasoning and suggest that gaze may function as a constraint rather than an enabler for capable MLLMs. Future work will explore adaptive gaze integration strategies, larger-scale evaluation on the full Main300 benchmark, and investigation of implicit gaze representations through visual attention mechanisms.

\bibliography{paper_acl_short/custom}
\bibliographystyle{plainnat}

\end{document}
